\documentclass[11pt,a4paper]{article}

\usepackage[utf8]{inputenc}
\usepackage[T1]{fontenc}
\usepackage[french]{babel}
\usepackage{xcolor}
\usepackage{titlesec}
\usepackage{geometry}
\usepackage{enumitem}
\usepackage{graphicx}
\usepackage{hyperref}
\usepackage{amssymb}

\geometry{margin=2.5cm}

% Couleurs Pangol1
\definecolor{primary}{RGB}{46, 125, 50}
\definecolor{secondary}{RGB}{80, 80, 80}

\titleformat{\section}{\color{primary}\normalfont\Large\bfseries}{\thesection}{1em}{}
\titleformat{\subsection}{\color{primary}\normalfont\large\bfseries}{\thesubsection}{1em}{}
\titleformat{\subsubsection}{\color{primary}\normalfont\normalsize\bfseries}{\thesubsubsection}{1em}{}

\title{
    \color{primary}\Huge\textbf{Manuel d'Utilisation}\\[0.5cm]
    \Large Logiciel de Traitement de Données \textbf{Pangol1}\\[0.2cm]
    \small Guide Complet : Importation, Analyse et Visualisation
}
\author{Documentation Technique Pangol1}
\date{\today}

\begin{document}

\maketitle

\begin{abstract}
\textbf{Pangol1} est une solution logicielle open source de haute performance spécialisée dans l'importation, le traitement et la visualisation de données structurées massives.
\end{abstract}

\tableofcontents
\newpage

% =========================
\section{Présentation de l'Interface}
% =========================

L'interface de Pangol1 est conçue pour l'analyse de gros volumes de données. Elle génère ensuite ses données sous forme de graphique définit via des filtres choisis par l'utilisateur

\subsection{Barre de Menu Supérieure}

\begin{itemize}
    \item \textbf{Fichier :} Ouvrir un fichier local, charger depuis une URL, fermer une table.
    \item \textbf{Affichage :} Modifier la disposition des fenêtres et de la mise en forme de l'application pour le confort de l'utilisateur
    \item \textbf{Aide :} Changer la langue et ouvrir la documentation.
\end{itemize}

\subsection{L'Explorateur de Données}
La partie gestion de données ou se trouve votre fichier chargé. Cette partie vous permet de trier ou de filtrer les données de votre fichier.
\begin{itemize}
    \item Pagination des données
    \item Filtres
    \item Gestion des colonnes
\end{itemize}

\subsection{Le Panneau d'Actions}
La partie d'affichage et de gestion des graphiques. Elle vous permet de :
\begin{itemize}
    \item Ajouter un graphique et personnaliser son affichage. 
    \item Ajouter une image
    \item Ajouter un texte
\end{itemize}

\subsection{Le Carnet de note}
Vous permet de visualiser vos donnée sous forme de graphique, d'image ou textuel.

% =========================
\section{Configuration de l'Application}
% =========================

\subsection{Fenêtre Paramètres}

\subsubsection{Général}
Conseillé à la première utilisation de l'application. Choississez votre langue ainsi que vos préférences de démarrage et de fermeture d'application.

\subsubsection{Table}
\begin{itemize}
    \item Sensibilité à la conversion de type correspond à l'exactitude des calculs que l'ordinateur met dans sa performance de conversion de type.
    \item Afficher ou non les numéros de lignes, les valeurs nulles.
    \item Ajuster les préférences pour les colonnes
    \item Caculer automatiquement les Statistiques des colonnes
\end{itemize}

\subsubsection{Apparence}

\begin{itemize}
    \item Thème clair / sombre / automatique
    \item Heures du mode sombre
    \item Activation FlatLaf
    \item Taille de police
    \item Famille de police
\end{itemize}

\subsubsection{Raccourcis}

\begin{itemize}
    \item Ctrl + O : Ouvrir
    \item Ctrl + W : Fermer
    \item Ctrl + F : Filtres
    \item Ctrl + , : Paramètres
    \item Ctrl + Shift + A : Tri croissant
    \item Ctrl + Shift + D : Tri décroissant
    \item Ctrl + Z : Annuler
    \item Ctrl + Y : Rétablir
    \item Vous pouvez aussi ajouter ou modifier vos propre raccourcis
\end{itemize}

% =========================
\section{Guide d'Utilisation : Importation et Traitement}
% =========================

\subsection{Importation d'un fichier CSV/Parquet/txt/html}

\begin{enumerate}
    \item Importer un fichier via "Fichier/Ouvrir" ou une URL via "fichier/Ouvrir une URL".
    \item Traitement automatique des données. Si erreur voir les conditions de lecture de format disponible (CSV/Parquet/txt/html).
\end{enumerate}

\subsection{Manipulation des Colonnes}

\begin{itemize}
    \item \textbf{Filtre:} Permet de créer des filtres sur les données chargées au sein d'une colonne.
    \item \textbf{Afficher/Masquer les colonnes:} Permet d'afficher ou de masquer les colonnes vides au sein des données chargées.
    \item \textbf{Clic droit sur une colonnes:} Ouvre en détail le panneau de gestion de donnée : 
    \begin{itemize}
    \item \textbf{Tri:} Ascendent/Descendent, soit du plus petit au plus grand.
    \item \textbf{Changer de type de données:} Transforme le type de donnée de la colonnes en un autre ( ex: entier -> chiffre à virgule). Certain type ne sont pas compatible au changement et la colonnes deviendra vide
    \item \textbf{Filtre:} Affecte des filtres personnalisés à la colonne selectionnée.
    \item \textbf{Statistiques:} Le taux de cardinalité, ou de valeurs nulles (nombres de ligne vide). Ainsi qu'un bouton résumé référencent le : min,max,moyenne,sommes,nombres,valeurs nulles,type, et type natif de la colonne
    \end{itemize}
\end{itemize}

% =========================
\section{Module d'Analyse et Statistiques}
% =========================

\subsection{Statistiques}
En bas de l'écran est écrit le nom du fichier charger, ainsi que le nombres de lignes des colonnes et le nombre de colonne chargé.
Un résumé des Statistiques d'une colonne est résumé en bas de l'application en fonction du type de donnée selectionnée au sein de la colonne :
\begin{itemize}
    \item \textbf{Nombre Entier:} Moyenne, minimum de la colonne, maximum de la colonne
    \item \textbf{Date:} date la plus récente, date la plus ancienne
    \item \textbf{Chaine de caractère:} sommes, moyenne, minimum de la colonne, maximum de la colonne
\end{itemize}

Il est aussi possible d'avoir plus de détail d'une colonne en se référent à la section 3.2 ci-dessus.

% =========================
\section{Visualisation Graphique}
% =========================
Le bouton \textbf{Ajouter un graphique:} vous permet d'ajouter et de personnaliser un graphique à l'application. Selectionner une colonne en cliquant dessus pour selectionner la colonnes à transformer en graphique

\subsection{Types}

\begin{itemize}
    \item \textbf{Types:} Barres, Camembert, Ligne, Nuages de Points, Radar, Aires
    \item \textbf{Titre:} Titre du graphique
\end{itemize}


\subsection{Configuration}

\begin{itemize}
    \item \textbf{Echelle:} Linéraire, logarithmique, racines carrée
    \item Nombre de graduation sur les axes
    \item Afficher/ne pas afficher les Axes. Et les renommer.
\end{itemize}

\subsection{Ajouter une Image/Texte}
Les boutons \textbf{Ajouter une image/Ajouter un bloc de texte}, permettent dans le carnet de note de rajouter un bloc de texte personnaliser, ou d'importer une image selectionner.

% =========================
\section{Carnet de Notes et Export}
% =========================

Le bouton \textbf{Ajouter au carnet de note} lors de la selection d'un résumé d'une donnée permet de l'ajouter à la fenêtres d'affichage (carnet de note) et de l'exporter.
Le bouton \textbf{d'export} s'affiche en haut du panneau d'affichage lorsqu'une données est affiché ( graphique/image/texte). Elle permet d'exporter le carnet de note au format : HTML, PDF ou markdown.
De lui donner un thème parmis : Lin, Encre, Sauge, Braise, Havre, Grès. De personnaliser le nom du fichier et enfin d'inclure ou non le titre d'un graphique.

% =========================
\section{Support Technique}
% =========================

Pour tout problème d'application,
contactez l'un des administrateur système \textbf{"Aide/A Propos"}.

\end{document}